\chapter{Rééchantillonnage}\label{ch:qm:resampling}

Une stratégie qui vend de la volatilité affiche un \omterm{def:qm:estimation:sharpe}{ratio de Sharpe} de 1.1 sur cinq ans, et le responsable du desk demande un intervalle à 95~\%
avant de doubler son capital. Rééchantillonner les cinq années jour par jour donne un intervalle allant de 0.17 à 2.25~: largement positif.
Rééchantillonner des blocs d'environ un mois, ce qui conserve ensemble les semaines calmes et les semaines violentes, donne $-0.25$ à $2.96$~: les données
ne permettent pas d'exclure que la stratégie n'ait aucun edge. Seul le second intervalle est honnête, et encore est-il trop étroit. Ce chapitre traite de
l'estimation de l'incertitude d'une statistique en la recalculant sur les données elles-mêmes~: le \omterm{def:qm:resampling:bootstrap}{bootstrap} et pourquoi il fonctionne, les variantes par blocs
que demandent les données dépendantes, les \omterm{def:qm:resampling:perm}{tests de permutation} et l'\omterm{def:qm:resampling:perm}{échangeabilité} qui les rend exacts, le \omterm{def:qm:resampling:jackknife}{jackknife}, plus ancien, et la
version \omterm{def:qm:resampling:bootstrap}{bootstrap} du test du chapitre~12 pour le meilleur parmi de nombreuses stratégies.

\section{Le bootstrap}

\begin{definition}[Fonction de répartition empirique, bootstrap]\label{def:qm:resampling:bootstrap}
La \emph{fonction de répartition empirique}\index{fonction de répartition empirique} de $X_1, \dots, X_n$ est $F_n(x) = \frac1n\#\{i: X_i
\le x\}$, la loi qui place une masse $1/n$ sur chaque observation. Le \emph{bootstrap}\index{bootstrap} estime la loi d'échantillonnage d'une
statistique $\hat\theta = T(X_1, \dots, X_n)$ par la loi de $\hat\theta^* = T(X_1^*, \dots, X_n^*)$, où les $X_i^*$ sont tirés
indépendamment de $F_n$, c'est-à-dire avec remise à partir des données~; en pratique en calculant $\hat\theta^*$ sur $B$ rééchantillons.
\end{definition}

Le principe est une substitution~: la loi inconnue $F$ est remplacée par $F_n$, et l'expérience d'échantillonnage qu'on ne peut pas répéter dans
le monde est rejouée sur ordinateur. L'\omterm{def:qm:estimation:estimator}{erreur type} \omterm{def:qm:resampling:bootstrap}{bootstrap} d'une moyenne est $s\sqrt{(n - 1)/n}/\sqrt n$, la formule usuelle~; son usage
est réservé aux statistiques dépourvues de formule, ou dont la formule repose sur des hypothèses auxquelles on ne se fie pas.

\begin{definition}[Intervalle bootstrap des percentiles, probabilité de couverture]\label{def:qm:resampling:interval}
L'\emph{intervalle bootstrap des percentiles}\index{intervalle bootstrap des percentiles} de niveau $1 - \alpha$ va du quantile $\alpha/2$ au quantile $1 -
\alpha/2$ des tirages \omterm{def:qm:resampling:bootstrap}{bootstrap} $\hat\theta^*$. La \emph{probabilité de couverture}\index{probabilité de couverture} d'une procédure d'intervalle
est la probabilité, sur des échantillons répétés tirés de la vraie loi, que son intervalle contienne le vrai paramètre.
\end{definition}

Un intervalle nominal à 95~\% est une affirmation sur la couverture, et celle-ci se mesure par simulation~: on tire de nombreuses trajectoires issues d'une loi connue
on construit l'intervalle sur chacune et l'on compte. C'est le test auquel chaque intervalle de ce chapitre est soumis.

Le \omterm{def:qm:resampling:bootstrap}{bootstrap} est convergent (la loi de $\sqrt n(\hat\theta^* - \hat\theta)$ conditionnellement aux données se rapproche de celle de $\sqrt n(\hat\theta -
\theta)$) pour les fonctionnelles régulières de $F$ à données i.i.d.~: moyennes, variances, corrélations, \omterm{def:qm:estimation:sharpe}{ratios de Sharpe}, coefficients de régression. Il échoue
lorsque la statistique dépend des valeurs extrêmes. Pour le maximum d'un échantillon uniforme, chaque rééchantillon manque le maximum de l'échantillon avec
probabilité $(1 - 1/n)^n \to e^{-1}$, de sorte que la loi \omterm{def:qm:resampling:bootstrap}{bootstrap} de $n(\max^* - \max)$ présente un atome de masse $1 - e^{-1} = 0.632$ en zéro,
alors que la vraie loi de $n(\theta - \max)$ est continue. Les quantiles de queue, les drawdowns et la plus forte perte partagent ce problème~: le
\omterm{def:qm:resampling:bootstrap}{bootstrap} ne peut pas rééchantillonner un krach que l'historique ne contient pas.

\section{Données dépendantes~: bootstrap par blocs et bootstrap stationnaire}

Rééchantillonner les jours indépendamment détruit tout ce qui relie un jour au suivant. Pour une stratégie dont le P\&L est autocorrélé, le
\omterm{def:qm:resampling:bootstrap}{bootstrap} i.i.d. estime alors la variance de la moyenne par $\hat\gamma(0)/n$ où la vérité est la variance de long terme sur $n$
(chapitre~11), et ses intervalles sont trop étroits.

\begin{definition}[Bootstrap par blocs mobiles, bootstrap stationnaire]\label{def:qm:resampling:blocks}
Le \emph{bootstrap par blocs mobiles}\index{bootstrap par blocs mobiles} (Künsch, 1989) construit un rééchantillon en concaténant des blocs de $b$
observations consécutives, $X_{j+1}, \dots, X_{j+b}$, dont les points de départ $j$ sont tirés uniformément (indices pris modulo $n$ dans la
version circulaire). Le \emph{bootstrap stationnaire}\index{bootstrap stationnaire} (Politis and Romano, 1994) tire les longueurs de blocs
indépendamment d'une loi géométrique de moyenne $b$~: chaque jour rééchantillonné prolonge le bloc en cours avec une probabilité $1 - 1/b$ et
commence un nouveau bloc à une position uniforme avec une probabilité $1/b$.
\end{definition}

\begin{proposition}[Le bootstrap stationnaire est stationnaire]\label{prop:qm:resampling:stationary}
Conditionnellement aux données, la série issue du \omterm{def:qm:resampling:blocks}{bootstrap stationnaire} $X_1^*, X_2^*, \dots$ est strictement stationnaire.
\end{proposition}

\begin{proof}
Soit $I_t$ l'indice de $X_t^*$. $I_1$ est uniforme sur $\{1, \dots, n\}$. Sachant $I_{t-1}$, $I_t$ vaut $I_{t-1} + 1$ (modulo $n$) avec
probabilité $1 - 1/b$ et uniforme avec probabilité $1/b$~; les deux applications envoient la loi uniforme sur elle-même, de sorte que chaque $I_t$ est uniforme. La
suite $(I_t)$ est une \omterm{def:qm:markov-chains-and-queues:ctmc}{chaîne de Markov} démarrée sous sa \omterm{def:qm:stochastic-differential-equations:kolmogorov}{loi stationnaire}, donc strictement stationnaire, et il en va de même pour $X_t^* = X_{I_t}$.
\end{proof}

Les blocs de longueur fixe du \omterm{def:qm:resampling:blocks}{bootstrap par blocs mobiles} rendent le rééchantillon non stationnaire~: les raccords entre blocs tombent à des
dates fixes. Les longueurs aléatoires éliminent cet artefact au prix d'une variance un peu plus grande, et les deux méthodes ont besoin d'une longueur de bloc. Trop courte, et
la dépendance est coupée~; trop longue, et les blocs distincts à recombiner sont peu nombreux. L'erreur quadratique moyenne de l'\omterm{def:qm:estimation:estimator}{estimateur} de la variance est
minimisée par une longueur de bloc croissant comme $n^{1/3}$, avec une constante qui dépend des autocovariances~; Politis et White (2004),
avec la correction de Patton, Politis et White (2009), l'estiment à partir des autocovariances empiriques au moyen d'une fenêtre de retards à sommet plat.
C'est la règle que met en œuvre \texttt{optimal\_block\_length}.

\begin{example}[Le vendeur de volatilité]\label{ex:qm:resampling:seller}
Chaque jour, la stratégie encaisse une prime et verse le carré du rendement du jour, $x_t = \mu + 1 - r_t^2$, avec $r_t = \sigma_t\varepsilon_t$ et
$\ln\sigma_t$ suit un AR(1) de coefficient 0.9 et d'écart type stationnaire 0.4, normalisé de sorte que $\E[\sigma_t^2] = 1$. Les rendements de marché
$r_t$ sont non corrélés, mais les rendements au carré ne le sont pas~: $\Cov(r_t^2, r_{t+k}^2) = e^{0.64 \cdot 0.9^k} - 1$, et la variance de long terme du P\&L
vaut 3.9 fois sa variance. La prime $\mu$ est calibrée pour un \omterm{def:qm:estimation:sharpe}{ratio de Sharpe} de population de 1.1. L'historique de cinq ans simulé à graine fixe
présente un \omterm{def:qm:estimation:sharpe}{ratio de Sharpe} de 1.11, une asymétrie de $-4.7$, un kurtosis de 34 et des autocorrélations de 0.13 à un jour, 0.10 à cinq jours,
$-0.02$ à vingt jours.
\end{example}

La règle automatique retient une longueur moyenne de bloc de 18.6 jours pour le \omterm{def:qm:resampling:blocks}{bootstrap stationnaire} sur cet historique (21.2 pour la version circulaire), et
les distributions \omterm{def:qm:resampling:bootstrap}{bootstrap} se séparent (\cref{fig:qm:resampling:dist})~: le \omterm{def:qm:resampling:bootstrap}{bootstrap} i.i.d. donne une \omterm{def:qm:estimation:estimator}{erreur type} de 0.53 et
l'intervalle $(0.17, 2.25)$~; le \omterm{def:qm:resampling:blocks}{bootstrap par blocs mobiles} 0.84 et $(-0.34, 2.93)$~; le \omterm{def:qm:resampling:blocks}{bootstrap stationnaire} 0.83 et $(-0.25, 2.96)$.
La \omterm{met:qm:estimation:delta}{méthode delta} du chapitre~11, avec l'asymétrie et le kurtosis empiriques, concorde avec le \omterm{def:qm:resampling:bootstrap}{bootstrap} i.i.d. (0.52)~: les deux supposent des
jours indépendants.

\begin{omfigure}
\begin{tikzpicture}
\begin{axis}[width=10.5cm,height=5cm,xlabel={ratio de Sharpe bootstrap (annualisé)},ylabel={densité},
  tick label style={font=\small},label style={font=\small},xmin=-1.5,xmax=4,ymin=0,ymax=0.9,grid=major,grid style={gray!25},
  legend columns=3,legend style={font=\footnotesize,at={(0.5,-0.32)},anchor=north,draw=none,fill=none,
  /tikz/every even column/.append style={column sep=8pt}},table/col sep=comma]
\addplot[omThm,thick,const plot mark mid,fill=omThm!15] table[x=x,y=iid]{figdata/methods/13-resampling/dist.csv};
\addplot[omDef,very thick,const plot mark mid] table[x=x,y=stationary]{figdata/methods/13-resampling/dist.csv};
\addplot[black,dashed,thin,sharp plot] coordinates {(1.106,0) (1.106,0.9)};
\legend{bootstrap i.i.d.,{stationnaire, blocs de 18.6 jours},valeur de l'échantillon}
\end{axis}
\end{tikzpicture}

\omcaption{Distributions \omterm{def:qm:resampling:bootstrap}{bootstrap} du \omterm{def:qm:estimation:sharpe}{ratio de Sharpe} sur cinq ans du vendeur de volatilité (valeur de l'échantillon 1.11, en tirets)~: rééchantillonnage indépendant des
jours (en grisé, \omterm{def:qm:estimation:estimator}{erreur type} 0.53) et rééchantillonnage de blocs aléatoires de longueur moyenne 18.6 jours (\omterm{def:qm:estimation:estimator}{erreur type} 0.83). Données~: le
tutoriel du chapitre, 4\,000 rééchantillons, à graine fixe.}\label{fig:qm:resampling:dist}
\end{omfigure}

Laquelle des méthodes est juste relève de la simulation. Sur 4\,000 trajectoires du même processus, le \omterm{def:qm:estimation:sharpe}{ratio de Sharpe} sur cinq ans a un
écart type de 1.10~: supérieur à toutes les estimations \omterm{def:qm:resampling:bootstrap}{bootstrap}. Moyennée sur soixante trajectoires, l'\omterm{def:qm:estimation:estimator}{erreur type} du \omterm{def:qm:resampling:blocks}{bootstrap stationnaire}
croît avec la longueur de bloc, de 0.59 (blocs d'un jour, soit le \omterm{def:qm:resampling:bootstrap}{bootstrap} i.i.d.) jusqu'à un palier d'environ 0.97 entre 40 et 90 jours, et
n'atteint jamais la vraie valeur (\cref{fig:qm:resampling:blocks}). Le déficit tient aux queues~: les trajectoires qui, par hasard, ne contiennent aucune semaine violente
ne peuvent pas en rééchantillonner une.

\begin{omfigure}
\begin{tikzpicture}
\begin{semilogxaxis}[width=10cm,height=5cm,xlabel={longueur moyenne de bloc (jours)},ylabel={erreur type du ratio de Sharpe},
  tick label style={font=\small},label style={font=\small},xmin=1,xmax=120,ymin=0,ymax=1.3,grid=major,grid style={gray!25},
  xtick={1,2,5,10,20,40,60,120},log ticks with fixed point,
  legend style={font=\footnotesize,at={(0.98,0.05)},anchor=south east,fill=white,draw=none},table/col sep=comma]
\addplot[omDef,very thick,mark=*,mark size=1.5pt] table[x=block,y=se]{figdata/methods/13-resampling/blocks.csv};
\addplot[black,dashed,thick] table[x=block,y=truth]{figdata/methods/13-resampling/blocks.csv};
\legend{bootstrap stationnaire (moyenne),écart type d'échantillonnage vrai}
\end{semilogxaxis}
\end{tikzpicture}

\omcaption{\omterm{def:qm:estimation:estimator}{Erreur type} du \omterm{def:qm:estimation:sharpe}{ratio de Sharpe} sur cinq ans par \omterm{def:qm:resampling:blocks}{bootstrap stationnaire}, en fonction de la longueur moyenne de bloc, moyennée sur soixante
historiques simulés du vendeur de volatilité, et l'écart type vrai (1.10, sur 4\,000 historiques). Les blocs d'un jour correspondent au
\omterm{def:qm:resampling:bootstrap}{bootstrap} iid (0.59)~; le plateau voisin de 0.97 est le maximum que retrouve une longueur de bloc quelconque. Données~: le tutoriel du chapitre,
graine fixée.}\label{fig:qm:resampling:blocks}
\end{omfigure}

Le verdict est donné par la couverture (la \cref{fig:qm:resampling:coverage}). Sur 400 historiques simulés, les intervalles iid à 95~\% nominal contiennent la valeur vraie
du \omterm{def:qm:estimation:sharpe}{ratio de Sharpe} de 1.1 dans 68.5~\% d'entre eux~; les intervalles par blocs mobiles et stationnaires, à longueur de bloc automatique, dans 82.75~\%~; les
intervalles stationnaires à blocs de 40 jours dans 84.0~\%. Supprimer le regroupement de volatilité (un coefficient AR nul) et l'intervalle iid couvre
91.5~\%, le déficit restant étant dû aux queues épaisses.

\begin{omfigure}
\begin{tikzpicture}
\begin{axis}[width=9cm,height=4.4cm,xbar,bar width=9pt,
  xlabel={part des intervalles à 95~\% contenant le ratio de Sharpe vrai (\%)},
  ytick={0,1,2,3},yticklabels={iid,{blocs mobiles, automatique},{stationnaire, automatique},{stationnaire, 40 jours}},
  tick label style={font=\small},label style={font=\small},
  xmin=0,xmax=100,ymin=-0.6,ymax=3.6,grid=major,grid style={gray!25},table/col sep=comma]
\addplot[fill=omDef!60,draw=omDef,nodes near coords,nodes near coords style={font=\footnotesize,anchor=east,xshift=-1pt}]
  table[x=coverage,y=k]{figdata/methods/13-resampling/coverage.csv};
\addplot[black,dashed,thick,sharp plot] coordinates {(95,-0.6) (95,3.6)};
\end{axis}
\end{tikzpicture}

\omcaption{Couverture des intervalles \omterm{def:qm:resampling:bootstrap}{bootstrap} à 95~\% nominaux pour le \omterm{def:qm:estimation:sharpe}{ratio de Sharpe} du vendeur de volatilité, sur 400 historiques simulés de cinq ans~:
iid 68.5~\%, \omterm{def:qm:resampling:bootstrap}{bootstraps} par blocs à longueur automatique 82.75~\%, stationnaire à blocs de 40 jours 84.0~\%~; la ligne en pointillés marque le niveau
nominal de 95~\%. Données~: le tutoriel du chapitre, graine fixée.}\label{fig:qm:resampling:coverage}
\end{omfigure}

\section{Tests de permutation}

\begin{definition}[Échangeabilité, test de permutation]\label{def:qm:resampling:perm}
Les variables aléatoires $Y_1, \dots, Y_n$ sont dites échangeables (\emph{échangeabilité}\index{echangeabilite@échangeabilité}) si la loi de $(Y_{\pi(1)}, \dots,
Y_{\pi(n)})$ est la même pour toute permutation $\pi$. Un \emph{test de permutation}\index{test de permutation} de l'\omterm{def:qm:testing-and-multiple-testing:test}{hypothèse nulle} selon laquelle $Y$ est
indépendant de $X$ compare une statistique $T(X, Y)$ à ses valeurs sur des permutations aléatoires de $Y$~; la $p$-valeur vaut $(1 + \#\{T(X,
Y_\pi) \ge T(X, Y)\})/(1 + P)$ pour $P$ permutations.
\end{definition}

\begin{proposition}[Exactitude sous l'échangeabilité]\label{prop:qm:resampling:exact}
Si, sous l'\omterm{def:qm:testing-and-multiple-testing:test}{hypothèse nulle}, $Y$ est échangeable et indépendant de $X$, la $p$-valeur de permutation vérifie $\P(p \le \alpha) \le \alpha$ pour
tout $\alpha$ et toute taille d'échantillon.
\end{proposition}

\begin{proof}
Sous l'\omterm{def:qm:testing-and-multiple-testing:test}{hypothèse nulle}, le $Y$ observé et les $P$ copies permutées sont échangeables conditionnellement à $X$, donc le rang de $T(X, Y)$ parmi les $P + 1$
valeurs est uniforme (égalités départagées en défaveur du rejet). La $p$-valeur est ce rang compté depuis le haut, divisé par $P + 1$, et $\P(\text{rang} \le
k) \le k/(P + 1)$.
\end{proof}

La condition est à elle seule tout le test. Le P\&L quotidien n'est pas échangeable lorsqu'il est dépendant dans le temps, et un signal persistant ne l'est pas non plus.
Prenons le P\&L du vendeur de volatilité et un signal construit sur un autre marché, indépendant~: sa moyenne des carrés de rendements sur les vingt derniers jours.
Rien ne les relie. Pourtant, sur 200 couples indépendants de ce type, le \omterm{def:qm:resampling:perm}{test de permutation} de leur corrélation rejette au seuil de 5~\% dans 21.5~\% d'entre eux,
car la permutation détruit la persistance qui rend grandes les corrélations fallacieuses entre deux séries lentes. Décaler une série
d'un décalage aléatoire (une rotation circulaire, qui conserve sa dépendance) rejette dans 2.5~\% des cas. Le même piège touche les tests d'un signal contre
des rendements qui se chevauchent, et le remède est le même~: rééchantillonner de façon à préserver ce que l'\omterm{def:qm:testing-and-multiple-testing:test}{hypothèse nulle} ne rompt pas.

\section{Le jackknife}

\begin{definition}[Jackknife]\label{def:qm:resampling:jackknife}
Le \emph{jackknife}\index{jackknife} recalcule une statistique en omettant tour à tour chaque observation, $\hat\theta_{(i)}$, et, avec $\bar\theta
= \frac1n\sum_i\hat\theta_{(i)}$ estime son biais par $(n - 1)(\bar\theta - \hat\theta)$ et son \omterm{def:qm:estimation:estimator}{erreur type} par
$\sqrt{\frac{n-1}n\sum_i(\hat\theta_{(i)} - \bar\theta)^2}$.
\end{definition}

\begin{proposition}[Correction de biais par jackknife]\label{prop:qm:resampling:jack}
Si $\E[\hat\theta_n] = \theta + a/n + O(n^{-2})$, l'estimation corrigée $\hat\theta - (n - 1)(\bar\theta - \hat\theta)$ a pour biais
$O(n^{-2})$.
\end{proposition}

\begin{proof}
On écrit $\E[\hat\theta_n] = \theta + a/n + c/n^2 + O(n^{-3})$. Chaque estimation obtenue en omettant une observation utilise $n - 1$ observations, donc $\E[\bar\theta] =
\theta + a/(n - 1) + c/(n - 1)^2 + O(n^{-3})$. L'estimation corrigée vaut $n\hat\theta - (n - 1)\bar\theta$, de moyenne $\theta + a - a +
c\bigl(1/n - 1/(n - 1)\bigr) + O(n^{-2}) = \theta + O(n^{-2})$~: les termes en $1/n$ s'annulent exactement.
\end{proof}

Pour la variance par substitution $\frac1n\sum(X_i - \bar X)^2$, dont le biais vaut $-\sigma^2/n$, la correction restitue l'\omterm{def:qm:estimation:estimator}{estimateur sans biais} $s^2$
exactement. Sur l'historique du vendeur de volatilité, l'\omterm{def:qm:estimation:estimator}{erreur type} \omterm{def:qm:resampling:jackknife}{jackknife} du \omterm{def:qm:estimation:sharpe}{ratio de Sharpe} vaut 0.53, soit celui du \omterm{def:qm:resampling:bootstrap}{bootstrap} iid, et
pour la même raison~: en omettant des jours isolés, il les traite comme indépendants. Le \omterm{def:qm:resampling:jackknife}{jackknife} par blocs, qui supprime des blocs, en est la version
dépendante. Le \omterm{def:qm:resampling:jackknife}{jackknife} est peu coûteux ($n$ recalculs, sans tirage aléatoire) et adapté au biais~; le \omterm{def:qm:resampling:bootstrap}{bootstrap} fournit des distributions entières.

\section{Bootstrap du Reality Check}

Le test de statistique maximale du chapitre~12 simulait l'\omterm{def:qm:testing-and-multiple-testing:test}{hypothèse nulle} de deux cents règles de momentum à partir d'une loi gaussienne dotée de leur corrélation estimée.
Le \omterm{def:qm:testing-and-multiple-testing:rc}{Reality Check} fait de même avec les données~: il rééchantillonne conjointement les jours pour les deux cents règles (de sorte que la corrélation entre règles soit
préservée), au moyen d'un \omterm{def:qm:resampling:blocks}{bootstrap stationnaire} (afin que toute dépendance dans le temps soit conservée), recentre la moyenne rééchantillonnée de chaque règle sur sa moyenne empirique (afin que
les rééchantillons vérifient l'\omterm{def:qm:testing-and-multiple-testing:test}{hypothèse nulle}), puis retient la statistique $t$ maximale de chaque rééchantillon.

Avec 2\,000 rééchantillons par \omterm{def:qm:resampling:blocks}{bootstrap stationnaire}, le P\&L des règles de momentum étant presque indépendant dans le temps (la longueur automatique de bloc est
1.3 jours), la $p$-valeur ajustée de la meilleure règle vaut 0.0275, contre 0.030 pour la simulation gaussienne, et la valeur critique à 5~\% du
maximum vaut 2.85 contre 2.92 (la \cref{fig:qm:resampling:rc}). L'approximation gaussienne suffisait ici~; avec des P\&L à queues épaisses ou
dépendants dans le temps, elle ne suffirait plus, et le \omterm{def:qm:resampling:bootstrap}{bootstrap} n'a besoin d'aucune modification.

\begin{omfigure}
\begin{tikzpicture}
\begin{semilogyaxis}[width=10cm,height=5cm,xlabel={seuil $c$},ylabel={$\P(\max t \ge c)$ sous l'hypothèse nulle},
  tick label style={font=\small},label style={font=\small},xmin=0,xmax=3.6,ymin=1e-3,ymax=1.5,grid=major,grid style={gray!25},
  legend style={font=\footnotesize,at={(0.03,0.05)},anchor=south west,fill=white,draw=none},table/col sep=comma]
\addplot[omDef,very thick] table[x=c,y=bootstrap]{figdata/methods/13-resampling/rc.csv};
\addplot[omThm,very thick,dashed] table[x=c,y=gaussian]{figdata/methods/13-resampling/rc.csv};
\addplot[black!50,thin,sharp plot] coordinates {(3.11,1e-3) (3.11,1.5)};
\legend{bootstrap stationnaire,gaussienne à corrélation estimée}
\end{semilogyaxis}
\end{tikzpicture}

\omcaption{Loi sous l'\omterm{def:qm:testing-and-multiple-testing:test}{hypothèse nulle} de la meilleure statistique $t$ parmi les deux cents règles de momentum du chapitre~12, estimée par 2\,000 rééchantillons par \omterm{def:qm:resampling:blocks}{bootstrap stationnaire}
de leur P\&L conjoint et par 200\,000 tirages gaussiens à corrélation estimée. Au 3.11 de la meilleure règle (trait vertical)~:
0.0275 et 0.030. Données~: le tutoriel du chapitre, graine fixée.}\label{fig:qm:resampling:rc}
\end{omfigure}

\section{Tutoriel~: l'intervalle trop étroit}\label{tut:qm:resampling:interval}

\begin{tutorial}
\textbf{Objectif.} Faire un \omterm{def:qm:resampling:bootstrap}{bootstrap} du \omterm{def:qm:estimation:sharpe}{ratio de Sharpe} du vendeur de volatilité de trois façons et mesurer la \omterm{def:qm:resampling:interval}{probabilité de couverture} de chaque intervalle.
\textbf{État final~:} les \cref{fig:qm:resampling:dist,fig:qm:resampling:blocks,fig:qm:resampling:coverage}~; couvertures de 68.5~\% et 82.75~\%.
\begin{tutsteps}
\item \textbf{Indices de rééchantillonnage} pour les blocs à longueur fixe et à longueur géométrique, générés pour tous les rééchantillons à la fois, afin qu'une matrice de P\&L à plusieurs colonnes
soit rééchantillonnée de façon cohérente.
\omcode{code/firm/resample/firm_resample.py}{29}{47}{Indices du bootstrap par blocs mobiles et du bootstrap stationnaire.}\label{lst:qm:resampling:indices}
\item \textbf{La longueur des blocs} à partir des autocovariances empiriques, avec une fenêtre à sommet plat.
\omcode{code/firm/resample/firm_resample.py}{55}{83}{Longueur moyenne automatique des blocs (Politis--White, corrigée).}\label{lst:qm:resampling:block}
\item \textbf{Exécuter} \texttt{problem()}, \texttt{coverage()} et \texttt{coverage(block=40)} dans \texttt{qm\_resampling.py}, puis
\texttt{fig\_resampling.py}.
\end{tutsteps}
\textbf{Pour aller plus loin.} Remplacer l'\omterm{def:qm:resampling:interval}{intervalle bootstrap des percentiles} par un \omterm{def:qm:resampling:bootstrap}{bootstrap} studentisé (rééchantillonner la statistique $t$, et non le ratio) et
en mesurer la \omterm{def:qm:resampling:interval}{probabilité de couverture}~; fixer le coefficient AR de la log-volatilité à 0.97 et constater que chaque intervalle se dégrade.
\end{tutorial}

\section{Construction~: le module de rééchantillonnage}\label{bld:qm:resampling:resample}

\begin{build}
\bfield{Rôle} Intervalles et $p$-valeurs pour les statistiques de la firme miniature qui n'ont pas de formule fiable~: \omterm{def:qm:estimation:sharpe}{ratios de Sharpe} de
stratégies à P\&L dépendants, drawdowns, le meilleur résultat d'une recherche.
\bfield{Interface} \texttt{iid\_indices}, \texttt{block\_indices}, \texttt{stationary\_indices} (chacune \texttt{(n, n\_boot, \dots, rng)}
renvoyant un tableau d'entiers de forme \texttt{(n\_boot, n)})~; \texttt{optimal\_block\_length(x)}~; \texttt{bootstrap(stat, x, idx)}~;
\texttt{percentile\_interval(draws, level)}~; \texttt{permutation\_test}, \texttt{circular\_shift\_test}~; \texttt{jackknife(stat, x)}.
\bfield{Règles} Des indices, et non des données rééchantillonnées, afin que chaque colonne d'une matrice de P\&L soit rééchantillonnée avec les mêmes jours~; un générateur est toujours
transmis en argument~; aucun intervalle n'est publié sans que sa \omterm{def:qm:resampling:interval}{probabilité de couverture} ait été mesurée sur un modèle des données.
\bfield{Tests d'acceptation} \texttt{code/firm/resample/tests/}~: longueurs de bloc géométriques de la moyenne demandée~; lois marginales uniformes à chaque
position d'un rééchantillon stationnaire~; l'\omterm{def:qm:estimation:estimator}{erreur type} \omterm{def:qm:resampling:bootstrap}{bootstrap} d'une moyenne~; la règle de longueur des blocs sur bruit blanc et sur un AR(1)~; la
taille exacte d'un \omterm{def:qm:resampling:perm}{test de permutation}~; le \omterm{def:qm:resampling:jackknife}{jackknife} reproduit $s^2$ exactement.
\bfield{Pour aller plus loin} Intervalles studentisés et intervalles accélérés corrigés du biais~; le \omterm{def:qm:resampling:jackknife}{jackknife} par blocs~; le test SPA de Hansen en complément de
\texttt{stationary\_indices} et \texttt{firm.multitest}.
\end{build}

\begin{omsources}
\item B.~Efron, «~Bootstrap methods: another look at the jackknife~», \emph{Annals of Statistics} 7, 1979.
\item H.~R.~Künsch, «~The jackknife and the bootstrap for general stationary observations~», \emph{Annals of Statistics} 17, 1989.
\item D.~N.~Politis et J.~P.~Romano, «~The stationary bootstrap~», \emph{Journal of the American Statistical Association} 89, 1994.
\item D.~N.~Politis et H.~White, «~Automatic block-length selection for the dependent bootstrap~», \emph{Econometric Reviews} 23, 2004~;
correction par A.~Patton, D.~N.~Politis et H.~White, \emph{Econometric Reviews} 28, 2009.
\item M.~H.~Quenouille, «~Approximate tests of correlation in time-series~», \emph{Journal of the Royal Statistical Society B} 11, 1949.
\end{omsources}

\section{Exercices}

\begin{exercise}[$\star$]\label{exo:qm:resampling:1}
Déterminer la probabilité qu'un jour donné soit absent d'un rééchantillon \omterm{def:qm:resampling:bootstrap}{bootstrap} iid de $n$ jours, puis la calculer pour $n = 1\,260$ et
dans la limite.
\end{exercise}

\begin{exercise}[$\star$]\label{exo:qm:resampling:2}
Dans un \omterm{def:qm:resampling:blocks}{bootstrap stationnaire} de longueur moyenne de bloc égale à 20 jours, déterminer la probabilité qu'un bloc dure plus de 60 jours.
\end{exercise}

\begin{exercise}[$\star$]\label{exo:qm:resampling:3}
Soit un échantillon de cinq observations. Dénombrer les rééchantillons \omterm{def:qm:resampling:bootstrap}{bootstrap} ordonnés, puis les rééchantillons distincts, vus comme des multiensembles.
\end{exercise}

\begin{exercise}[$\star\star$]\label{exo:qm:resampling:4}
Montrer que la correction de biais par \omterm{def:qm:resampling:jackknife}{jackknife} de la variance plug-in $\frac1n\sum(X_i - \bar X)^2$ donne $s^2 = \frac1{n-1}\sum(X_i - \bar X)^2$.
\end{exercise}

\begin{exercise}[$\star\star$]\label{exo:qm:resampling:5}
Pour un AR(1) de coefficient 0.5 et $n = 1\,260$, calculer la longueur moyenne optimale de bloc du \omterm{def:qm:resampling:blocks}{bootstrap stationnaire} à partir de $b =
(G/g(0))^{2/3}n^{1/3}$, avec $G = \sum_k|k|\gamma(k)$ et $g(0) = \sum_k\gamma(k)$.
\end{exercise}

\begin{exercise}[$\star\star$]\label{exo:qm:resampling:6}
Pour $n = 100$ observations uniformes, déterminer la probabilité qu'un rééchantillon \omterm{def:qm:resampling:bootstrap}{bootstrap} contienne le maximum de l'échantillon. Commenter
ce résultat pour le \omterm{def:qm:resampling:bootstrap}{bootstrap} du maximum.
\end{exercise}

\begin{exercise}[$\star\star\star$]\label{exo:qm:resampling:7}
\emph{Programmation.} Sur 200 couples indépendants de P\&L d'un vendeur de volatilité et d'un signal de variance glissante issu d'un autre marché, mesurer à quelle fréquence
un \omterm{def:qm:resampling:perm}{test de permutation} et un test par décalage circulaire de leur corrélation rejettent au seuil de 5~\%.
\end{exercise}

\begin{exercise}[$\star\star\star$]\label{exo:qm:resampling:8}
\emph{Trouver la faille.} «~Nous avons appliqué le \omterm{def:qm:resampling:bootstrap}{bootstrap} à nos cinq années de P\&L quotidien dix mille fois. Le 5\up{e}~percentile des \omterm{def:qm:estimation:sharpe}{ratios de Sharpe} rééchantillonnés
est de 0.4, donc avec un niveau de confiance de 95~\%, le \omterm{def:qm:estimation:sharpe}{ratio de Sharpe} de la stratégie dépasse 0.4.~»
\end{exercise}

\section{Problème~: l'intervalle trop étroit}

\begin{problem}[{Problème du week-end --- intervalle sur le ratio de Sharpe d'un vendeur de volatilité}]\label{pb:qm:resampling:1}
Une stratégie encaisse chaque jour une prime et verse le carré du rendement de marché du jour~; la log-volatilité du marché suit un processus AR(1) de
coefficient 0.9 et d'écart-type 0.4. Le \omterm{def:qm:estimation:sharpe}{ratio de Sharpe} de la population vaut 1.1, et les cinq années de P\&L quotidien sont
l'historique simulé du chapitre, à graine fixée.

\medskip\noindent\textbf{Partie I --- L'historique.}
\begin{enumerate}
\item Calculer le \omterm{def:qm:estimation:sharpe}{ratio de Sharpe} empirique, le coefficient d'asymétrie et le kurtosis.
\item Calculer les autocorrélations à un, cinq et vingt jours.
\item Expliquer pourquoi le P\&L est autocorrélé alors que les rendements du marché ne le sont pas.
\item Calculer la variance de long terme du P\&L dans la population, en multiples de sa variance.
\item Calculer l'\omterm{def:qm:estimation:estimator}{erreur type} donnée par la \omterm{met:qm:estimation:delta}{méthode delta}, en utilisant l'asymétrie et le kurtosis empiriques.
\end{enumerate}

\medskip\noindent\textbf{Partie II --- Trois intervalles.}
\begin{enumerate}[resume]
\item Calculer l'intervalle \omterm{def:qm:resampling:bootstrap}{bootstrap} iid et son \omterm{def:qm:estimation:estimator}{erreur type}.
\item Déterminer les longueurs de bloc retenues par la règle automatique.
\item Calculer l'intervalle du \omterm{def:qm:resampling:blocks}{bootstrap par blocs mobiles} et son \omterm{def:qm:estimation:estimator}{erreur type}.
\item Calculer l'intervalle du \omterm{def:qm:resampling:blocks}{bootstrap stationnaire} et son \omterm{def:qm:estimation:estimator}{erreur type}.
\item Calculer l'\omterm{def:qm:estimation:estimator}{erreur type} et le biais donnés par le \omterm{def:qm:resampling:jackknife}{jackknife}.
\end{enumerate}

\medskip\noindent\textbf{Partie III --- \omterm{def:qm:resampling:interval}{Probabilité de couverture}.}
\begin{enumerate}[resume]
\item Calculer l'écart-type réel du \omterm{def:qm:estimation:sharpe}{ratio de Sharpe} sur cinq ans.
\item Déterminer à quelle fréquence les intervalles iid contiennent 1.1, sur 400 historiques.
\item Faire de même pour les intervalles stationnaires, avec la longueur automatique puis avec des blocs de 40 jours.
\item Expliquer pourquoi aucune longueur de bloc n'atteint 95~\%.
\item Expliquer ce que devient l'intervalle iid en l'absence de regroupement de volatilité.
\end{enumerate}

\medskip\noindent\textbf{Partie IV --- Jugement.}
\begin{enumerate}[resume]
\item Quel intervalle donnez-vous au responsable du desk~?
\item Faut-il doubler le capital au vu de ces éléments~?
\item Qu'ajouteriez-vous à l'analyse avant de décider~?
\item Énoncer le \emph{résultat nommé}~: la \omterm{def:qm:resampling:interval}{probabilité de couverture} des intervalles nominaux à 95~\% iid et \omterm{def:qm:resampling:blocks}{bootstrap stationnaire}.
\item En une phrase~: que doit préserver un schéma de rééchantillonnage~?
\end{enumerate}
\end{problem}

\section{Questions d'entretien}

\begin{interviewq}[$\star$ \iqroles{chercheur, ingénieur ML}]\label{iq:qm:resampling:1}
Que fait le \omterm{def:qm:resampling:bootstrap}{bootstrap}, et quand échoue-t-il~?
\end{interviewq}

\begin{interviewq}[$\star\star$ \iqroles{chercheur, trader}]\label{iq:qm:resampling:2}
Pourquoi ne pas faire un \omterm{def:qm:resampling:bootstrap}{bootstrap} indépendant des rendements quotidiens d'une stratégie pour obtenir un intervalle sur son \omterm{def:qm:estimation:sharpe}{ratio de Sharpe}~?
\end{interviewq}

\begin{interviewq}[$\star\star$ \iqroles{chercheur}]\label{iq:qm:resampling:3}
Comment choisissez-vous la longueur de bloc d'un \omterm{def:qm:resampling:bootstrap}{bootstrap} par blocs~?
\end{interviewq}

\begin{interviewq}[$\star\star$ \iqroles{chercheur, ingénieur ML}]\label{iq:qm:resampling:4}
Vous testez si un signal prédit les rendements en mélangeant le signal. Quand ce test est-il trompeur~?
\end{interviewq}

\begin{interviewq}[$\star\star$ \iqroles{chercheur}]\label{iq:qm:resampling:5}
\omterm{def:qm:resampling:jackknife}{Jackknife} ou \omterm{def:qm:resampling:bootstrap}{bootstrap}~: que vous apporte chacun~?
\end{interviewq}

\begin{interviewq}[$\star\star\star$ \iqroles{chercheur}]\label{iq:qm:resampling:6}
Comment calculeriez-vous une $p$-valeur pour le meilleur de deux cents backtests qui respecte à la fois la corrélation entre eux et
la dépendance dans le temps~?
\end{interviewq}
